\begin{table}[htbp]
  \centering
  \caption{長距離条件（C2）と幅狭条件（C3）の第4ブロック平均および改善量の比較（各条件10名）．差はC2からC3を引いた値，95\%信頼区間は個別・未補正のWelch区間である．対応のない両側Welchの\(t\)検定を用い，6比較にHolm補正を適用した．}
  \label{tab:practice-c2-c3-between}
  \small
  \setlength{\tabcolsep}{3pt}
  \begin{tabular}{lrrrrrr}
    \hline
    比較項目 & C2平均 & C3平均 & 差 [95\%信頼区間] & \(t(df)\) & \(p\) & \(p_{\mathrm{Holm}}\) \\
    \hline
    \(MT\) 第4ブロック平均（ms） & 1074.39 & 1245.03 & \(-170.64\;[-519.50,\,178.22]\) & \(-1.04\,(15.41)\) & 0.3143 & 1.0000 \\
    \(ER\) 第4ブロック平均（\%） & 12.60 & 24.03 & \(-11.43\;[-20.96,\,-1.90]\) & \(-2.57\,(13.89)\) & 0.0222 & 0.1332 \\
    \(TP_{\mathrm{e}}\) 第4ブロック平均（\(\mathrm{s}^{-1}\)） & 51.11 & 39.41 & \(11.70\;[-7.43,\,30.82]\) & \(1.29\,(16.87)\) & 0.2140 & 1.0000 \\
    \hline
    \(MT\) 短縮量（ms） & 334.25 & 250.25 & \(84.01\;[-171.08,\,339.09]\) & \(0.69\,(17.51)\) & 0.4972 & 1.0000 \\
    \(ER\) 低下量（ポイント） & 3.10 & 5.57 & \(-2.47\;[-14.25,\,9.31]\) & \(-0.46\,(11.00)\) & 0.6536 & 1.0000 \\
    \(TP_{\mathrm{e}}\) 増加量（\(\mathrm{s}^{-1}\)） & 17.03 & 13.07 & \(3.96\;[-10.68,\,18.59]\) & \(0.57\,(18.00)\) & 0.5769 & 1.0000 \\
    \hline
  \end{tabular}
\end{table}
